\section{临朝、改号与复唐}
\label{sec:empress-wu-zhou}

655年高宗废王皇后、立武后，666年二人封禅泰山，显示后宫、礼仪与中央人事已不能截然分开。680年废李贤、改立李显，683年高宗死后武后以皇太后临朝，继承秩序遂成为公开的政治问题。正史后来往往以性别与正统的语言裁决这些变化；可确认的则是皇后、太后能够利用诏令、宗室、官僚与礼仪进入权力中心，而其控制也始终需要这些中介。%
\footnote{立后、封禅、废立与临朝见 ST-655-WU-EMPRESS、ST-666-TAISHAN、ST-680-LI-XIAN、ST-683-GAOZONG-DEATH，依据《旧唐书》《新唐书》高宗、中宗纪及《资治通鉴》。宫闱动机和仪式言辞带有强烈的后出叙事。}

684年徐敬业在扬州起兵而败，690年武后改唐为周、自立为帝。\eventanchor{ST-690-ZHOU-FOUNDATION}{武周建国（690）}这不是两件能脱离地方军政的宫廷新闻：扬州的反叛显示新的权力安排须经过区域道路与兵源的检验；受命请愿与改号则把官僚、宗教性符瑞和礼制共同编入新的合法性语言。请愿文本能说明朝廷如何展示同意，不能证明所有参与者的意愿或各地的服从密度。%
\footnote{徐敬业起兵与建周见 ST-684-LI-JINGYE、ST-690-ZHOU-FOUNDATION，依据《旧唐书·则天皇后纪》《新唐书·则天皇后纪》及《资治通鉴》。反叛规模、符瑞和请愿程序需警惕胜者文本的编排。}

\begin{debate}{改周是否等于全部政治关系断裂？}
复唐正统语言常把690年写成必须纠正的例外；请愿、符瑞与受命文本则服务于新政权的合法性展示。它们能说明权力怎样被表述，不能量出各地官员和社会群体的认同，更不能将唐、周之间的行政延续抹去。
\end{debate}

705年张柬之等迫武则天退位、恢复唐号，707年李重俊又兵败，说明复唐并未立即消除宫廷集团的竞争。年号的恢复重新接上李唐宗室与既有官僚的记忆，却不等于所有人事、财政和地方关系回到高宗以前。武周的历史更适合被理解为继承、礼仪和行政资源被重新组合的一段时间，而非可以由一个“篡”字穷尽的例外。%
\footnote{复唐与李重俊失败见 ST-705-TANG-RESTORATION、ST-707-LI-CHONGJUN，依据《旧唐书·中宗本纪》《旧唐书·张柬之传》《新唐书》相关纪传及《资治通鉴》。政变执行细节和太子部众构成并不完整。}

\subsection{交叉阅读窗口三：改号的文书与幽州台上的孤声}

武周的名号不能只由一部正史裁定。《旧唐书》《新唐书》的则天纪与《资治通鉴》共同保存临朝、建周和复唐的关键年序，却都在复唐或北宋的正统语言中组织因果；《唐会要》中的礼制、选举条目和武周年号墓志，则可让我们看见诏令、官衔与纪年怎样实际进入书写。碑刻和墓志能够钉住某些人的任官与家世，佛教造像题记能够显示符瑞语言的流通，却不能证明请愿者真心一致，亦不能量出地方服从。%
\footnote{可互校《旧唐书·则天皇后纪》《新唐书·则天皇后纪》《资治通鉴》高宗、中宗纪与《唐会要》相关礼制、选举条目；武周年号墓志和造像题记适于核验个案，不宜外推为全国政治态度。}

陈子昂《登幽州台歌》说：“前不见古人，后不见来者。念天地之悠悠，独怆然而涕下。”这是一首古体短歌，通常联系到万岁通天元年（696）北方行军的挫折，具体作地与作年仍有讨论。其抒情力度来自把一位随军官员的失意压缩为无前无后的时间感；它并非武周反对派的民意陈述，更不能据此断定边军或河北社会如何看待女皇。钱穆关于政治制度何以取得承认的长线问题，在此不宜化为性别道德判断，而应追问改号、年号、诏命、荐举和边防如何被共同使用。诗保存的是一个识字官员的姿态，墓志和制度文书保存的是不同的制度痕迹；二者的距离，正是武周政治不能被一声“篡夺”抹平的证据限度。
